\documentclass[11pt]{article}
\usepackage{amsmath,amssymb,amsthm,booktabs,hyperref}
\newtheorem{theorem}{Theorem}
\newtheorem{lemma}[theorem]{Lemma}
\newtheorem{proposition}[theorem]{Proposition}
\newtheorem{corollary}[theorem]{Corollary}
\newcommand{\sq}{\sqcup}
\title{Small halves and three intrinsic centre strata}
\author{Human Compression project}
\date{}
\begin{document}
\maketitle

\paragraph{Scope and attribution.}
Throughout, a target graph means a connected simple nonbipartite quintic
graph of order $32$, girth four and diameter eight. We use the complete
classification relative to a fixed eight-geodesic in~\cite{rows}, its
intrinsic centre theorem~\cite{centre}, and the retained empty and
singleton middle results~\cite{middle}. The original cage problem and
its computational census belong to Cambie, Goedgebeur, Jooken and
Van den Eede~\cite{original}. Those results are not attributed to the
present note.

The new work here is the ordinary identification of the halves in three
centre strata. Two strata have six ordinary graphs each, by the finite
proofs below. The third has a product formula in four finite half counts;
we do not evaluate those counts. Standard ingredients are incidence
margins, independent subsets, graph distance, and preservation of the
centre under isomorphisms. The centre-four and centre-path arguments
were suggested in the project discussion and checked independently in
written assessments~\cite{assessments}; we give complete proofs rather
than treating those assessments as enumeration receipts.

No subject Source, graph enumeration, TeX compilation or Lean verification
was executed in preparing this note. An analytic orbit count proved here
does not assert an independent computational census, completion of the
original manuscript, or formal verification.

\section{The interface and the retained reduction}

Write $C(G)$ for the centre of $G$. The radius is four. A false-twin
class means an independent set of vertices having the same full open
neighborhood. The intrinsic centre theorem proves that $G-C(G)$ has
exactly two connected components, each of order at most sixteen.
The centre has at most two full-graph false-twin classes. An outside
vertex meets all or none of each such class. Its boundary color records
which classes it meets.

An isomorphism of boundary-colored components preserves their ordinary
edges and their boundary colors; it need not preserve a selected endpoint,
row, side, or fixed geodesic. The intrinsic gluing theorem of~\cite{centre}
states that the weighted centre graph together with its two components,
colored in this way, is a complete ordinary graph invariant. If centre
classes have equal size, their coordinate permutation must be applied
simultaneously to both components. We will only use a single centre
class or two classes of unequal sizes, so there is no such ambiguity here.

For clarity we state the remaining imported local notation. Fix endpoints
$p_0,p_8$ at distance eight and a midpoint $p_4$. Their neighborhoods
$A,B,C$ form an independent set of size fifteen. The induced residual
graph on the other seventeen vertices has exactly five edges, and the
three named vertices are isolated in that residual graph. The other
fourteen vertices split into two endpoint sides and a middle set $M$.
Every side vertex has positive incidence with its endpoint neighborhood;
both sides have size at least five. An endpoint row has four neighbors
in its side, forming an independent subset, and its fifth neighbor is
its endpoint. We use $q=|M|$ only to identify the imported cases, not as
an ordinary invariant.

In the nonisolated middle classification, the complete centre table
implies the following facts.
\begin{enumerate}
\item A centre consisting of one false-twin class of size four occurs
only in the retained all-isolated case $q=3$, or in the internal side
edge row with $q=4$, $j=1$.
\item A centre consisting of one false-twin class of size three occurs
only in the retained all-isolated case $q=2$, the no-middle-edge rows
$q=3$ of middle degree $d\in\{1,2,3,4\}$, or $Q_1,Q_2$.
\item A centre inducing $P_3$ occurs precisely in $Q_4,Q_7$. Its two
false-twin classes have sizes one and two, and the component orders
after deletion are thirteen and sixteen.
\end{enumerate}
These assertions combine~\cite{rows,centre} with the retained result
that the all-isolated middle cases have centre
$\{p_4\}\sq M$~\cite{middle}. The empty middle case is impossible and
the singleton case has centre size two. They therefore add none of the
three centre strata considered here.

Here and below a graph displayed with too few vertices is padded with
isolates to the stated side size. The notation $P_k$ means the path on
$k$ vertices, and $K_{1,k}$ is the star with $k$ leaves.

\begin{lemma}[A side edge inequality]\label{edgeineq}
Suppose a half has an endpoint $p$, five independent endpoint-row
vertices $R$, a side graph $D$, and $b$ independent port vertices $P$.
Every endpoint row has four neighbors in $D$, every port has an independent
neighbor subset in $D$, and every side vertex has total degree five.
Put
\[
 \alpha(v)=|N(v)\cap R|,\qquad
 \beta(v)=|N(v)\cap P|,\qquad d(v)=d_D(v).
\]
Then $\alpha(v)+\beta(v)+d(v)=5$. For every edge $uv$ of $D$,
\begin{equation}\label{edgeequation}
 d(u)+d(v)\geq5-b.
\end{equation}
\end{lemma}
\begin{proof}
An endpoint row cannot contain both endpoints of $uv$, so
$\alpha(u)+\alpha(v)\leq5$. A port cannot meet both endpoints, so
$\beta(u)+\beta(v)\leq b$. Adding the two side degree equations gives
$10\leq5+b+d(u)+d(v)$, as required.
\end{proof}

\section{A four-vertex centre forming one twin class}

Let the centre be a single false-twin class $X$ of size four. Removing
it gives two components of order fourteen. Across the two components
there are five boundary vertices, each adjacent to all of $X$ and hence
of internal degree one. We call them ports. The other vertices have
internal degree five. The two components have three and two ports,
respectively. This follows from the retained all-isolated $q=3$ forms
and the $q=4,j=1$ row in the imported centre table. We now classify both
components without retaining their endpoints as part of the structure.

The parameters in the all-isolated case can also be read off directly.
The side sizes sum to eleven and are at least five, hence are five and
six. A side of size $n$ with $e$ internal edges and $b$ singleton ports
has balance $b=5n-20-2e$. On the five-vertex side, two disjoint edges
would prevent an independent four-subset, and two adjacent edges would
force their common vertex into no endpoint row. Positive columns
therefore give $e\leq1$. Both port parts are nonempty; the five-vertex
side cannot have $e=0$, since that would use all five ports. Thus its
parameters are $(e,b)=(1,3)$, and the six-vertex side has $(e,b)=(4,2)$.

\begin{proposition}\label{leftfour}
There are exactly two possible ordinary three-port components. Each has
an endpoint $p$, five endpoint rows $R$, a five-vertex side
$D=K_2\sq3K_1$, and three ports of degree one. All ports attach to
endpoints of the side edge. Up to interchanging those endpoints, their
multiplicities are $(3,0)$ or $(2,1)$.
\end{proposition}
\begin{proof}
Both imported cases give this side. In the retained all-isolated case,
the five side vertices have four neighbors in each endpoint row, and
the side has one residual edge. In the $q=4,j=1$ row this is the side
with the internal edge and its three singleton side ports. The local
form is therefore the same in both cases.

Write the edge as $xy$. Every independent four-subset of a five-vertex
side omits $x$ or $y$. Each of its three isolates thus meets all five
rows, leaving no degree for a port. Let $\beta(x)+\beta(y)=3$. The
numbers of rows omitting $x,y$ must be $1+\beta(x),1+\beta(y)$,
since $\alpha(x)=4-\beta(x)$ and $\alpha(y)=4-\beta(y)$.
These numbers sum to five and determine all row multiplicities.
Up to $x\leftrightarrow y$, the two possible port splits are exactly
$(3,0)$ and $(2,1)$; both have nonnegative row multiplicities and positive
side columns.

Permuting equal-role ports and rows gives an isomorphism between any
two records with the same split. Forgetting $p$ cannot merge the two
splits: the port set is intrinsically the set of degree-one vertices,
and the number of distinct neighbors of the ports is one in the first
case and two in the second. Thus these are two ordinary unrooted
components, not merely two endpoint-marked records.
\end{proof}

\begin{proposition}\label{rightfour}
There are exactly three possible ordinary two-port components. In each,
the two ports have internal degree one. Their mutual distances are
respectively three, four and five; this distinguishes all three
components intrinsically.
\end{proposition}
\begin{proof}
In the retained $q=3$ case this component has an endpoint $p$, five
rows $R$, a six-vertex side $D$ with four edges, and two singleton
ports. Its endpoint columns are positive. In the $q=4,j=1$ case,
include its degree-four middle vertex in the six-vertex side. This
gives exactly the same format, now allowing a zero endpoint column.
Thus in both cases $D$ is triangle-free, has four edges, and
\begin{equation}\label{fourmargin}
 \alpha(v)+\beta(v)+d(v)=5,\qquad \alpha(v)\geq0,
 \qquad \sum_v\beta(v)=2.
\end{equation}
Every row is an independent four-subset. We classify the larger domain
defined by these conditions, then identify its cases with the imported
forms; this also proves sufficiency.

The triangle-free four-edge graphs on six vertices are
\[
 P_5\sq K_1,\quad K_{1,4}\sq K_1,\quad T\sq K_1,\quad
 C_4\sq2K_1,\quad P_4\sq K_2,\quad
 K_{1,3}\sq K_2,\quad 2P_3,
\]
where $T$ is the five-vertex tree with degree sequence $(3,2,1,1,1)$.
Indeed a cyclic graph must be $C_4$ plus isolates; a connected forest
with four edges is one of the three trees on five vertices; and the
remaining forest edge partitions are $3+1$ and $2+2$ as displayed.
Lemma~\ref{edgeineq}, with $b=2$, eliminates the two cases containing
an isolated edge. In $P_5\sq K_1$ the sole independent four-subset is
the three alternating path vertices together with the isolate. Thus
both other path vertices have $\alpha=0$, while their degrees are two
and $\beta\leq2$, contrary to~\eqref{fourmargin}. The same argument
excludes $2P_3$: its sole independent four-subset consists of the four
leaves, leaving the two degree-two centres with zero columns.

For $T\sq K_1$, let $c$ be the degree-three vertex, let $a,b$ be its
two leaf neighbors, and let the remaining branch be $cde$; let $f$ be
the isolate. No independent four-subset contains $c$, since only $e,f$
are nonneighbors of $c$. Hence $\alpha(c)=0$ and $\beta(c)=2$.
Every endpoint row avoiding $c$ must contain $a,b,f$ and one of $d,e$.
In particular $\alpha(a)=\alpha(b)=5$, inconsistent with their internal
degree one. This excludes $T$.

For $D=C_4\sq2K_1$, each independent four-subset consists of an
opposite cycle pair and both isolates. The isolates consequently have
$\alpha=5$ and $\beta=0$. If one opposite pair occurs $t$ times among
the five rows, its vertices have $\alpha=t$ and the other pair has
$\alpha=5-t$. Their port degrees are respectively $3-t$ and $t-2$.
Nonnegativity forces $t=2$ or $3$. Interchanging the opposite pairs
identifies these alternatives. Exactly one port meets each vertex of
one opposite pair, and no other vertex meets a port. This determines
one ordinary component. The two port neighbors are distance two in
the cycle, so the ports themselves are distance four.

For $D=K_{1,4}\sq K_1$, let $u$ be the centre and $f$ the isolate.
No independent four-subset contains $u$, so $\alpha(u)=0$ and
$\beta(u)=1$. Exactly one port meets $u$. The other port meets either
a star leaf or $f$, up to an automorphism of $D$. Conversely both
choices are feasible. The five rows are four-subsets of the five
vertices other than $u$; the row-omission multiplicity at $v$ is
$d(v)+\beta(v)$. Those multiplicities are nonnegative and sum to five.
For a leaf-port choice they are two at that leaf, one at the other
three leaves, and zero at $f$. For an isolate-port choice they are
one at each of the five vertices. The columns are positive except
for $u$ and every margin in~\eqref{fourmargin} is satisfied.

The two resulting ordinary components have port distances three and
five. In the first, the path is port--$u$--leaf--port. In the second,
$u$ and $f$ are at distance three through a leaf and an endpoint row:
they are not adjacent and have no common neighbor. Adding the two
port edges gives distance five. There are no further row choices,
since all omission multiplicities are forced. This proves exactly
three components and distinguishes them without a marked endpoint.
\end{proof}

\begin{lemma}\label{fourconverse}
Every pair of components in Propositions~\ref{leftfour} and
\ref{rightfour}, joined by attaching all five ports to each of four
independent new vertices $X$, is a target graph with centre exactly $X$.
\end{lemma}
\begin{proof}
All port degrees become five, and each new vertex has degree five.
The other degrees were already five. The graph is connected. Each
component is triangle-free by its independent endpoint rows and its
displayed triangle-free side graph; its ports are leaves, so the gluing
adds no triangle. Any two distinct rows on either side share a side
vertex and their endpoint, giving a four-cycle. The girth is four.

We verify the diameter and the exact centre directly. A port is at
distance one from $X$, and a side vertex meeting a port is at distance
two. On the three-port side the side edge has at least one port-supported
endpoint. Its other endpoint is therefore within three of $X$; the
three isolates meet all rows. Every row contains an edge endpoint,
so is within four of $X$. A row containing a port-supported endpoint
exists, and puts the side isolates within four as well.

On the cycle side, the two port-supported cycle vertices are opposite.
The other cycle vertices are within three of $X$. Each row contains
one opposite pair and the two isolates, and is within four of $X$;
the isolates also meet a row containing port support. On a star side,
the star centre is port-supported and every leaf is within three of
$X$. Every row contains at least three leaves and is within four of
$X$. The isolate is port-supported or meets every row; in the latter
case some row contains the other port's leaf neighbor and puts it
within four. Thus all side vertices and all endpoint rows are within
four in every case.

Each endpoint has distance exactly four from $X$. A path of length
four uses a row, a port-supported side vertex and a port; such a row
exists in every displayed form. There is no shorter path because the
endpoint meets only rows, each row meets only its endpoint and side,
and no row meets a port. Removing $X$ separates the endpoints, and
the false twins in $X$ have identical distances to either component.
Every path between endpoints consequently has length at least eight,
while a path through $X$ has length eight. All vertices are within four
of any fixed $x\in X$, so the diameter is exactly eight.

More precisely, if $v$ is in one component and $p$ is the endpoint of
the other, the separator identity of~\cite{middle} gives
$d(v,p)=d(v,x)+4>4$. Thus no outside vertex is central; all vertices
of $X$ are central and $C(G)=X$.

Finally the retained residual graph has the one left edge and the four
right edges. It is not $K_{1,5}$ because these edges lie in two distinct
components. The bipartiteness criterion in~\cite{rows} proves that the
graph is nonbipartite.
\end{proof}

\begin{theorem}\label{sixfour}
The target graphs whose centre is a single false-twin class of size
four form exactly six ordinary isomorphism classes.
\end{theorem}
\begin{proof}
The imported centre table places every such graph in the two local
formats classified above. Propositions~\ref{leftfour} and
\ref{rightfour} give two and three possible ordinary components,
and Lemma~\ref{fourconverse} realizes every pair. The centre is intrinsic.
Although the components both have order fourteen, their three and two
internal degree-one ports distinguish them. An ordinary isomorphism
therefore preserves the component roles and their unrooted types.
Conversely a pair of component isomorphisms extends over any permutation
of $X$, since every port is attached to all of $X$. There is no port
alignment to choose. The count is consequently $2\cdot3=6$.
\end{proof}

The cycle component belongs to the retained all-isolated middle case.
The two star components have a zero endpoint column and are supplied
by the $q=4,j=1$ internal side edge row. They are combined here by
their common intrinsic centre rather than by their different middle
descriptions.

\section{A centre inducing a three-vertex path}

Let $C(G)$ induce $P_3$. Write $C_1$ for its singleton middle class
and $C_2$ for its two false-twin leaves. Their unequal sizes distinguish
the classes intrinsically. A vertex meeting $C_1$ only will be called
an $A$-port; one meeting both vertices of $C_2$ only will be called a
$B$-port. Their internal degrees after deletion are four and three,
respectively. All other component vertices have internal degree five.
The two components have orders thirteen and sixteen.

\begin{proposition}\label{thirteen}
There is exactly one boundary-colored ordinary thirteen-vertex component.
It consists of an endpoint $p$, five endpoint rows $R$, five independent
side vertices $S$, one $A$-port $u$ and one $B$-port $a$. The edge $ua$
is present. Their side neighborhoods are disjoint sets of sizes three
and two, partitioning $S$. The row--side graph is $K_{5,5}$ with a
perfect matching deleted.
\end{proposition}
\begin{proof}
In row $Q_4$, the original degree-three middle vertex is $u$ and the
sole side port is $a$. Their edge is prescribed by that row. The side
arities are three and two, and their supports are disjoint because a
common neighbor would give a triangle. Hence every side vertex meets
one of the two ports. Its degree equation requires exactly four endpoint
row neighbors. Each endpoint row omits one side vertex, so the five
omissions use each side vertex once. This is exactly the stated crown
graph.

Any two choices of the $3+2$ partition are mapped to one another by a
permutation of $S$. Matching the row that omits each vertex extends this
map to an isomorphism. This proves uniqueness for $Q_4$. In $Q_7$, the
sole original side port is the $A$-port and the degree-two middle vertex
is the $B$-port. They are adjacent and have side arities three and two,
again with complementary supports. The same argument gives precisely
the same colored ordinary component.
\end{proof}

\begin{proposition}\label{sixteen}
There are exactly six boundary-colored ordinary sixteen-vertex components.
They are indexed by an ordered pair $(\lambda,\mu)$, where
\[
 \lambda\in\{(2),(1,1)\},\qquad
 \mu\in\{(3),(2,1),(1,1,1)\}.
\]
\end{proposition}
\begin{proof}
In $Q_4$ this component has an endpoint $p$, five endpoint rows $R$,
five independent side vertices $S$, two $A$-ports $u_1,u_2$, and three
$B$-ports $v_1,v_2,v_3$. The induced port graph is $K_{2,3}$. Each
port has exactly one additional neighbor, lying in $S$, because its
internal degree is four or three. Opposite-color ports are adjacent,
so they cannot share a side neighbor. Ports of the same color may do
so. Thus the equality multiplicities of the two $A$-port neighbors
form a partition $\lambda$ of two, and those of the three $B$-port
neighbors form a partition $\mu$ of three. Their supports are disjoint.
Every displayed pair is possible on five side vertices, since it uses
at most $2+3=5$ distinct vertices.

For $s\in S$ let $\delta(s)$ be the total number of ports using it.
Degree five forces $\alpha(s)=5-\delta(s)$ endpoint rows at that
vertex. Every row omits exactly one side vertex. Therefore the row
omission multiset is forced: it contains $s$ with multiplicity
$\delta(s)$. Its total size is five. The maximum value is three,
since supports of the two colors are disjoint; hence all endpoint
columns are positive. This proves feasibility and leaves no additional
row choice. The selected rows are independent because $S$ is independent.

These partitions are intrinsic under ordinary colored isomorphism.
The ports are exactly the vertices of internal degrees four and three,
with the two colors fixed by their centre attachments. They induce the
intrinsic $K_{2,3}$. Every port has a unique neighbor outside that port
set. The equality relation among those neighbors, separately for each
color, recovers $\lambda$ and $\mu$. Different partition pairs thus
cannot be isomorphic, even when endpoints and row labels are forgotten.

Conversely equal partition pairs permit permutations within each port
color and a bijection of the distinct used side vertices. Extend this
by a bijection of unused side vertices. The map preserves $\delta$,
so matching rows with their omitted vertex extends it to an isomorphism
of the full components. This proves that the partition pair is complete.

In $Q_7$, the two original opposite-side ports become the $A$-ports
and the three degree-one middle vertices become the $B$-ports. Their
induced graph is again $K_{2,3}$, and each has one side neighbor.
The same two partitions and forced omissions therefore give all its
components. In particular $Q_7$ adds no further ordinary component type.
\end{proof}

\begin{theorem}\label{sixpath}
The target graphs whose centre induces $P_3$ form exactly six ordinary
isomorphism classes.
\end{theorem}
\begin{proof}
The intrinsic centre table restricts this stratum to $Q_4,Q_7$.
Propositions~\ref{thirteen} and~\ref{sixteen} classify their components.
Every one of the six pairs is realized by $Q_4$: take the unique
thirteen-vertex component, any of the six sixteen-vertex components,
join every $A$-port to $C_1$, every $B$-port to both vertices of $C_2$,
and join $C_1$ to both vertices of $C_2$.

For an explicit identification with the imported construction, choose
one vertex $x\in C_2$ as the old midpoint, write the other as $h$, and
write $C_1=\{w\}$. The old middle set is
$\{h,u,u_1,u_2\}$, the old five midpoint-neighbors are
$\{w,a,v_1,v_2,v_3\}$, and the side middle degrees are $(3)$ and
$(1,1)$. The middle--port incidences are exactly $Q_4$. The thirteen
side has complementary supports, and on the sixteen side the two
opposite port colors have disjoint supports. The forced omissions
give positive endpoint columns and independent rows. Thus all local
hypotheses of the complete construction theorem~\cite{rows} hold.
It supplies simplicity, degree five, girth four, nonbipartiteness and
diameter eight; the intrinsic centre theorem gives exactly this $P_3$.

An ordinary isomorphism must preserve the centre, its unequal class
weights and its unequal component orders. It therefore preserves the
two boundary colors and the partition pair of the sixteen-component.
Conversely colored component isomorphisms extend over the centre;
there is no matching or alignment of ports, and the two leaves of
$C_2$ may be interchanged freely. Hence the six partition pairs give
exactly six ordinary graph classes. The $Q_7$ descriptions merely
change the provenance of centre and port roles, and give no additional
class.
\end{proof}

\section{A three-vertex centre forming one twin class}

This section gives a reusable finite domain rather than a new numerical
enumeration. For $b\in\{1,2,3,4\}$ define $\mathcal H_b$ as follows.
A record contains an endpoint $p$, five endpoint rows $R$, a six-vertex
side graph $D$ with $5-b$ edges, and $b$ ports. The endpoint meets all
five rows. Each row meets an independent four-subset of $D$, and each
port meets an independent two-subset of $D$. There are no further edges.
Require $D$ to be triangle-free and impose
\begin{equation}\label{hbmargin}
 \alpha(v)+\beta(v)+d(v)=5,\qquad\alpha(v)\geq0
 \quad(v\in V(D)).
\end{equation}
The count $h_b$ means the number of ordinary unrooted half graphs
represented by these records, with the ports retained intrinsically
as their degree-two vertices. It does not mean the number of records,
endpoint choices or incidence matrices. All the other $12$ vertices
have internal degree five, so each half has order $12+b$.

\begin{lemma}\label{zerob}
A record of $\mathcal H_b$ has at most one zero endpoint column.
If $\alpha(u)=0$, then $D$ is the star $K_{1,5-b}$ plus isolates,
with centre $u$, and every port contains $u$ in its two-subset.
All other side columns are positive. Every half in $\mathcal H_b$
is connected and triangle-free.
\end{lemma}
\begin{proof}
Put $e=5-b$. A zero column has $d(u)+\beta(u)=5$ and
$\beta(u)\leq b$, so $d(u)\geq5-b=e$. Since $D$ has exactly $e$
edges, equality holds: all its edges are incident with $u$, and
$\beta(u)=b$. Thus every port meets $u$, and its other neighbor lies
outside the star leaves by independence.

For $e\geq2$ the star has only one vertex of degree at least $e$.
For $e=1$ its two endpoints would both have to lie in every port if
both columns were zero, violating independence of a port's two-subset.
There is consequently at most one zero column.

The positive side vertices meet endpoint rows and connect to $p$.
Each port meets two distinct side vertices, at least one of which is
positive, so connects to $p$. A zero side vertex meets all ports and
also connects to $p$. This proves connectivity. No triangle can use
the endpoint, since its row neighbors are independent; no triangle
can use a row or port, since its side subset is independent; and $D$
itself is triangle-free. These exhaust the possible triangles.
\end{proof}

\begin{lemma}\label{dbtypes}
The possible side graph types in $\mathcal H_b$ lie in the following
list. A star with at least three leaves in the list has a zero centre
column. The listed types are subject to the incidence equations; the
list does not count their feasible ordinary half graphs.
\[
\begin{array}{c|c|l}
 b&e(D)&D\\\hline
1&4&C_4\sq2K_1\ \text{or }K_{1,4}\sq K_1\\
2&3&P_4\sq2K_1\ \text{or }K_{1,3}\sq2K_1\\
3&2&P_3\sq3K_1\ \text{or }2K_2\sq2K_1\\
4&1&K_2\sq4K_1
\end{array}
\]
In every feasible record the port support
$Z=\{v:\beta(v)>0\}$ is a vertex cover of $D$.
\end{lemma}
\begin{proof}
Lemma~\ref{edgeineq} gives $d(u)+d(v)\geq e(D)$ on each side edge.
For $e=4$, the seven four-edge types on six vertices were listed in
Proposition~\ref{rightfour}. Each of $P_5,T,P_4\sq K_2$,
$K_{1,3}\sq K_2$, and $2P_3$ has an edge whose endpoint degree sum
is less than four. Only the four-cycle and the four-leaf star survive.
For $e=3$ the four forest types are $P_4,K_{1,3},P_3\sq K_2,3K_2$;
the latter two have an isolated edge with degree sum two, and are
excluded. The two-edge and one-edge lists are exhaustive without
further restriction. A star centre of degree at least three cannot
belong to an independent four-subset on six vertices: at most its two
nonneighbors can join it in the three-leaf case, and at most its one
nonneighbor in the four-leaf case. Its column is therefore zero.

For a zero-column record, Lemma~\ref{zerob} makes the star centre
port-supported, so $Z$ covers every edge. Suppose instead all columns
are positive. For $b\geq2$, every side edge satisfies
$d(u)+d(v)\leq e(D)+1\leq4$. If both endpoints were outside $Z$,
the degree equations would give
$\alpha(u)+\alpha(v)=10-d(u)-d(v)\geq6$, impossible across five
independent rows. For $b=1$ the only positive type is $C_4\sq2K_1$,
and the same contradiction holds with endpoint degree sum four.
Thus $Z$ is always a vertex cover.
\end{proof}

\begin{lemma}\label{hbradius}
Join every port of a half in $\mathcal H_b$ to one new hub $x$.
Every vertex of the half is within four of $x$, and its endpoint
$p$ has distance exactly four from $x$.
\end{lemma}
\begin{proof}
Every port is at distance one and every vertex of $Z$ is at distance
two. A side vertex outside $Z$ which has a side edge has a neighbor
in $Z$ by Lemma~\ref{dbtypes}, and hence is within three. A side vertex
outside $Z$ with no side edge has $\alpha=5$ and meets all five rows.
At least one member of $Z$ has positive endpoint incidence: there are
at least two members of $Z$ because a port has two neighbors, and
Lemma~\ref{zerob} allows at most one zero column. A row containing
such a vertex is within three of $x$ and puts every side isolate
outside $Z$ within four.

A row meeting $Z$ is within three. If a row avoids $Z$, its four
side vertices exhaust the complement of $Z$, because $|Z|\geq2$ and
the side has size six. Consequently $|Z|=2$. Every port has that
same independent two-subset $Z$, and the complementary four-subset
is independent as the row. Every side edge must therefore run between
the two subsets. There is at least one such edge since $e(D)\geq1$.
Its endpoint in the row is within three of $x$, putting that row
within four. This treats all rows and side vertices.

The positive column in $Z$ supplies a path
$p$--row--side--port--$x$ of length four. There is no shorter path:
$p$ meets only rows, rows meet only $p$ and the side, and neither rows
nor $p$ meet a port. Thus $d(p,x)=4$.
\end{proof}

\begin{theorem}\label{threecentre}
Let $h_b$ be the ordinary unrooted half counts just defined. The target
graphs whose centre is one false-twin class of size three are exactly
the gluings of
\[
 \mathcal H_1\times\mathcal H_4
 \quad\text{and}\quad
 \mathcal H_2\times\mathcal H_3
\]
by attaching all five ports to all three vertices of an independent
centre $X$. Their ordinary graph count is
\begin{equation}\label{threecount}
 h_1h_4+h_2h_3.
\end{equation}
This formula contains no numerical evaluation of the $h_b$.
\end{theorem}
\begin{proof}
First consider necessity in the retained all-isolated case $q=2$.
The two side sizes sum to twelve. On a side of size $n$ with $e$
internal edges and $b$ ports, the port arity is two and the degree
balance is $2b=5n-20-2e$. Thus $n$ is even. Since both sides have
size at least five, each has size six. The balance gives $e=5-b$,
and the two port counts sum to five. This is a pair from the displayed
domains with all endpoint columns positive.

In a no-middle-edge $q=3$ row, take its degree-$d$ middle vertex on
the five-vertex side and include it in that side's six-vertex graph
$D$. Its $d$ side edges form a star and all $5-d$ own-side ports meet
it as well as one ordinary side vertex. The resulting component is
in $\mathcal H_{5-d}$ with one zero column. The opposite component
is in $\mathcal H_d$ with positive columns. In $Q_1,Q_2$, perform
the same inclusion on both sides. Their middle degrees are $(4,1)$
or $(3,2)$, giving the two displayed domain pairs with one zero column
in each half. The intrinsic centre table supplies no other one-class
three-centre case. Hence every graph in the stated stratum occurs.

Conversely take any pair from either displayed product. By
Lemma~\ref{zerob} its halves are connected and triangle-free. Add
three independent vertices $X$ and join every port to all three.
Each old degree-two port now has degree five, each new vertex meets
the five ports, and every other vertex already had degree five.
There are $3+(12+b)+(12+(5-b))=32$ vertices. There are no triangles
through $X$, since ports are independent within and across the two
halves. Distinct endpoint rows share side vertices, as two four-subsets
of a six-set intersect, giving a four-cycle through their endpoint.
Thus the graph is connected, simple, quintic and of girth exactly four.

Lemma~\ref{hbradius} puts every vertex within four of any fixed
$x\in X$, and puts both endpoints exactly four from $X$. Deleting
$X$ separates their connected halves, so their distance is eight.
The diameter is consequently eight. The separator identity gives
$d(v,p)=d(v,x)+4>4$ for any outside vertex $v$ and the opposite endpoint
$p$. Hence the centre is exactly the false-twin class $X$.

Using the endpoint row sets as $A,C$ and the five ports as $B$ gives
the retained independent fifteen-set. Its residual five edges are
the two disjoint side graphs, with $b$ and $5-b$ complementary edge
counts, both nonzero. They cannot form $K_{1,5}$. The retained
bipartiteness criterion therefore proves nonbipartiteness. This also
shows directly why zero columns may be permitted: they change which
vertices belong to the original middle set, without changing the
intrinsic centre or its two components.

Finally, the ports of a half are intrinsically its degree-two vertices.
The two component orders are $(13,16)$ or $(14,15)$, so neither the
component roles within a product nor the two products can merge under
ordinary isomorphism. Every isomorphism preserves the intrinsic centre
and induces component isomorphisms. Conversely any two component
isomorphisms extend across $X$, since all port attachments are full.
There are no gluing alignment choices. The two products and their
disjoint sum give~\eqref{threecount}.
\end{proof}

The symbols $h_b$ require ordinary identification across all endpoint
choices. Counting incidence records, summing the imported $q$ rows,
or declaring their descriptions disjoint would not evaluate them.
In contrast, Theorems~\ref{sixfour} and~\ref{sixpath} do give complete
ordinary analytic counts: their component invariants prove both
exhaustiveness and nonidentification directly.

\paragraph{Verification boundary.}
The proofs above are finite written mathematics using the explicitly
imported reduction and centre theorems. Independent review of this
edition is a separate artifact. Computational checks of these particular
half classifications and any Lean verification must be reported by their
own actual receipts; neither is asserted here. The full F009 ordinary
census still requires the other intrinsic centre strata and their
appropriate reviewed identification, or a separately complete union.

\begin{thebibliography}{9}
\bibitem{original} Cambie, Goedgebeur, Jooken and Van den Eede,
\emph{On the order-diameter ratio of girth-diameter cages},
arXiv:2511.21144v3, Appendix B, Table 6.
\bibitem{middle} Human Compression project,
\emph{The empty and singleton middle cases for quintic graphs},
frozen source SHA256
\texttt{76a5c09c2e3f52536c022572bd95ba611392ad38527c0f2ee40c1e886e9d9eb2},
with independent review SHA256
\texttt{026a34a5c0d80143005551a931a6eb5d83c899f29a048ed52aa72e075c5b6e63}.
\bibitem{rows} Human Compression project,
\emph{Nonisolated middle vertices in quintic graphs of diameter eight},
frozen source SHA256
\texttt{522af65bc2bb25d0d09ad68b1c047f2806832230e8452d388ead3483d84f6384},
with independent written review SHA256
\texttt{984216b5809678a74d06dee620592f0fc8f6511cace1ff1c3b9007072bcd496d}.
\bibitem{centre} Human Compression project,
\emph{The intrinsic centre and finite gluing interfaces},
frozen source SHA256
\texttt{25987bb9eb742cbc1b8e779408dec9823a99e3a4704e3eadb4f3d6735aa73a3d},
with independent written review SHA256
\texttt{82474f352b89e48a81194aac84b0eb43b8157262075e54ae5595a3a098a5131c}.
\bibitem{assessments} Separate project reviewer,
written assessments of the centre-four and centre-$P_3$ half arguments.
The latter is \emph{The intrinsic centre-$P_3$ stratum: independent
written derivation}, source SHA256
\texttt{b8dcda470c3d1bcdad53125159d66d2005c9fbc06568d33096d01454b1213a17}.
The present proofs include all their required arguments and do not
replace their own edition-level review by those assessments.
\end{thebibliography}
\end{document}
